% !TEX root = ../notes_sovereignai.tex
% Chapter 6: Next Steps. A curated roadmap of where to take the rig next.
%%%%%%%%%%%%%%%%%%%%%%%%%%%%%%%%%%%%%%%%%%%%%%%%%%%%%%%%%%%%%%%%%%%%%%

\index{roadmap}\index{multimodal}\index{OCR}\index{image generation}
\chapter{Next Steps}
\printchapterversion{06_nextsteps}
\newthought{Where the Rig Goes Next.}
\vspace{1.5em}

% ==============================================================================================
% THE WHY
% ==============================================================================================

\section{The Why}\index{extensibility}

\newthought{The stack is complete enough to live on}: a model on hardware you own, reachable
everywhere, answering by browser, terminal, voice, and from your own notes. This
last chapter is a map rather than a build, a list of where to take the rig when you
want more, each step a pointer with the tool and the reason. Every one of them is a
new client of the same engine or a new tenant on the same GPU, so none requires
starting over.

% ==============================================================================================
% THE ROADMAP
% ==============================================================================================

\section{The Roadmap}\index{vision}\index{fine-tuning}

\newthought{The capabilities worth adding next}, roughly in order of payoff for a
household:

\begin{itemize}
    \item Vision and OCR. A multimodal model such as Qwen3-VL\index{Qwen3}, pulled into Ollama
    like any other, lets the engine read images, screenshots, diagrams, and scanned
    or handwritten pages, turning a photo of a letter or a receipt into text and
    answers\cite{ollamaVision}.

    \item Documents and retrieval\index{knowledge base}. Beyond the Obsidian vault, point Open WebUI's
    knowledge feature at your PDFs and folders; it embeds and searches them so the
    family chat can answer from manuals, contracts, and records, all indexed on the
    rig\cite{openwebuiRAG}.

    \item Image generation\index{image generation}. ComfyUI\index{ComfyUI} on the same GPU adds text-to-image with SDXL\index{SDXL} or
    Flux\index{Flux}, a node-based pipeline that is faster and lighter on memory than the
    older web UIs and shares the card with the language engine\cite{comfyuiRepo}.

    \item The editor\index{coding agent}. Continue or Cline in VS Code points at the same Open WebUI API
    key from Chapter~3, so inline completion and chat in your editor run on the rig,
    the same engine your terminal already uses.

    \item Real concurrency. When the household load grows, move the serving layer to
    vLLM\index{vLLM}, whose continuous batching\index{continuous batching} serves many people at once from one
    card. Extra cards join as a replica each, or as a split by layer with pipeline
    parallelism\index{pipeline parallelism}, both of which leave the single lane alone where tensor
    parallelism\index{tensor parallelism} would not. It speaks the same API, so the clients above do not
    change\cite{vllmParallelism}.

    \item More ears. Add a Wyoming satellite\index{voice satellite} per room, custom wake words\index{wake word}, and
    streaming text-to-speech so replies begin before the sentence is finished; the
    Chapter~4 pipeline scales by adding microphones, not rebuilding.

    \item Health and hygiene. Watch GPU temperature and memory\index{monitoring}, back up the Docker
    volumes\index{backup} that hold accounts and history, keep models current with periodic pulls,
    and hold the cards near their undervolt\index{undervolting} so the rig runs cool and quiet for years.

    \item Your own domain. When a general model is not enough, fine-tune\index{fine-tuning} a small one
    or train a LoRA adapter\index{LoRA} on the rig against your own data, the deepest form of a
    model that is truly yours.
\end{itemize}

% ==============================================================================================
% CLOSING
% ==============================================================================================

\newthought{What you now own} is a machine in a room and an engine you understand
well enough to extend. Whatever you add from here runs on that hardware and keeps
working when the network does not.

\marginnote{\newthought{feedback}}
